\documentclass[11pt]{article}
\usepackage[margin=1in]{geometry}
\usepackage{amsmath,amssymb,amsthm}

\newtheorem*{theorem}{Theorem 4:2}
\newtheorem{lemma}{Lemma}

\title{A Formal Proof of Theorem 4:2 of Hestenes--Stiefel (1952)\\
\large Finite Termination of the Method of Conjugate Directions}
\date{}

\begin{document}
\maketitle

\section*{Setting}

Let $A \in \mathbb{R}^{n\times n}$ be symmetric positive definite, $k \in \mathbb{R}^n$,
and let $h$ be the unique solution of $Ah = k$. Write $(x,y) = x^{\top}y$.

The cd-method starts from $x_0$, sets $r_0 = k - Ax_0$, and for $i = 0, 1, \dots$ computes
\begin{align}
a_i &= \frac{(p_i, r_i)}{(p_i, Ap_i)}, \tag{4:1a}\\
x_{i+1} &= x_i + a_i p_i, \tag{4:1b}\\
r_{i+1} &= r_i - a_i A p_i, \tag{4:1c}
\end{align}
where the directions $p_0, p_1, \dots, p_{n-1}$ are nonzero and satisfy
\begin{equation}
(p_{i+1}, A p_j) = 0 \qquad (j = 0, 1, \dots, i). \tag{4:2}
\end{equation}
Such directions can always be chosen: if $p_0, \dots, p_i$ have been chosen with $i+1 < n$,
the set of vectors $A$-orthogonal to all of them is the orthogonal complement of
$\operatorname{span}\{Ap_0, \dots, Ap_i\}$, which has dimension at least $n - i - 1 \ge 1$.

From Theorem 4:1 we use
\begin{align}
(p_i, Ap_j) &= 0 \qquad (i \neq j), \tag{4:3a}\\
a_i &= \frac{(p_i, r_0)}{(p_i, Ap_i)}. \tag{4:4}
\end{align}

For $q = 0, 1, \dots, n$ define the subspaces
\[
V_q = \operatorname{span}\{p_0, \dots, p_{q-1}\}, \qquad V_0 = \{0\},
\]
and the initial error $y_0 = h - x_0$.

\begin{theorem}
Let $m$ be the smallest integer $q \in \{0, 1, \dots, n\}$ such that $y_0 \in V_q$. Then
\begin{enumerate}
\item[\textup{(i)}] $m$ is well defined and $m \le n$;
\item[\textup{(ii)}] $x_m = h$;
\item[\textup{(iii)}] $x_q \neq h$ for $q < m$.
\end{enumerate}
In particular, the cd-method is an $m$-step method with $m \le n$.
\end{theorem}

\section*{Proof}

\begin{lemma}\label{lem:indep}
The vectors $p_0, \dots, p_{n-1}$ are linearly independent. Consequently
$\dim V_q = q$ for $q = 0, \dots, n$, and $V_n = \mathbb{R}^n$.
\end{lemma}

\begin{proof}
Suppose $\sum_{i=0}^{n-1} c_i p_i = 0$. Fix $j$ and take the inner product with $Ap_j$.
By (4:3a) all terms with $i \neq j$ vanish, leaving $c_j (p_j, Ap_j) = 0$.
Since $p_j \neq 0$ and $A$ is positive definite, $(p_j, Ap_j) > 0$, so $c_j = 0$.
As $j$ was arbitrary, the vectors are independent. Hence $V_q$ is spanned by $q$
independent vectors, so $\dim V_q = q$, and $V_n$ is an $n$-dimensional subspace of
$\mathbb{R}^n$, i.e.\ $V_n = \mathbb{R}^n$.
\end{proof}

\begin{lemma}\label{lem:iterate}
For $q = 0, 1, \dots, n$,
\[
x_q = x_0 + \sum_{i=0}^{q-1} a_i p_i .
\]
In particular, $x_q - x_0 \in V_q$.
\end{lemma}

\begin{proof}
By induction on $q$. For $q = 0$ the sum is empty. If the formula holds for $q$, then
by (4:1b), $x_{q+1} = x_q + a_q p_q = x_0 + \sum_{i=0}^{q} a_i p_i$.
\end{proof}

\paragraph{Proof of (i).}
The set $S = \{ q \in \{0, \dots, n\} : y_0 \in V_q \}$ contains $n$, because
$V_n = \mathbb{R}^n$ by Lemma~\ref{lem:indep}. Thus $S$ is a nonempty finite set of
integers, so its minimum $m$ exists and $m \le n$. \hfill$\square$

\paragraph{Proof of (ii).}
If $m = 0$, then $y_0 \in V_0 = \{0\}$, so $x_0 = h$ and there is nothing to prove.

Let $m \ge 1$. Since $y_0 \in V_m$, there are scalars $\alpha_0, \dots, \alpha_{m-1}$ with
\begin{equation}
y_0 = \sum_{j=0}^{m-1} \alpha_j p_j. \tag{1}
\end{equation}
Because $Ah = k$,
\begin{equation}
r_0 = k - Ax_0 = A(h - x_0) = Ay_0 = \sum_{j=0}^{m-1} \alpha_j A p_j. \tag{2}
\end{equation}
Fix $i \in \{0, \dots, m-1\}$ and take the inner product of (2) with $p_i$. By (4:3a)
only the term $j = i$ survives:
\[
(p_i, r_0) = \alpha_i (p_i, Ap_i).
\]
Since $(p_i, Ap_i) > 0$, comparison with (4:4) gives
\begin{equation}
\alpha_i = \frac{(p_i, r_0)}{(p_i, Ap_i)} = a_i \qquad (i = 0, \dots, m-1). \tag{3}
\end{equation}
By Lemma~\ref{lem:iterate}, (3) and (1),
\[
x_m = x_0 + \sum_{i=0}^{m-1} a_i p_i = x_0 + \sum_{i=0}^{m-1} \alpha_i p_i = x_0 + y_0 = h.
\]
\hfill$\square$

\paragraph{Proof of (iii).}
Let $q < m$ and suppose $x_q = h$. By Lemma~\ref{lem:iterate},
$y_0 = h - x_0 = x_q - x_0 \in V_q$, so $q \in S$ with $q < m = \min S$,
a contradiction. Hence $x_q \neq h$. \hfill$\square$

\medskip
This completes the proof of Theorem 4:2. \hfill$\blacksquare$

\end{document}
